\documentclass[11pt,a4paper]{article}
\usepackage[utf8]{inputenc}
\usepackage[T1]{fontenc}
\usepackage[polish]{babel}
\usepackage[margin=2.5cm]{geometry}
\usepackage{array}
\usepackage{tabularx}
\usepackage{booktabs}

\setlength{\parindent}{0pt}
\setlength{\parskip}{0.6em}

\begin{document}

\noindent
\begin{tabularx}{\textwidth}{|p{0.13\textwidth}|X|p{0.30\textwidth}|}
\hline
\textit{grupa 6} & \textit{25.11.2026} & \textit{} \\[1.2em]
\hline
\multicolumn{3}{|l|}{\textit{4 - NMAP}} \\[1.2em]
\hline
\multicolumn{2}{|l|}{\textit{Filip Będkowski, Danylo Drygo}} & \textit{} \\[1.2em]
\hline
\end{tabularx}

\vspace{1em}

\section*{Wstęp teoretyczny}

\subsection*{Rodzina protokołów TCP/IP}

TCP/IP jest rodziną współpracujących protokołów, która stanowi podstawę
komunikacji w Internecie i wielu sieciach lokalnych. Umożliwia łączenie sieci
zbudowanych z użyciem różnych technologii warstwy łącza danych. Działanie
protokołów TCP/IP oraz zasady ich współpracy opisują dokumenty RFC
(\textit{Request For Comments}).

Podstawowym protokołem warstwy internetowej jest IP. W nagłówku pakietu IP
znajdują się między innymi adres źródłowy, adres docelowy i informacje
potrzebne do przekazania datagramu. Dane protokołów wyższej warstwy, takich
jak TCP lub UDP, są przenoszone jako ładunek pakietu IP. Routery przekazują
pakiety między sieciami na podstawie adresu docelowego i tablic routingu.
W praktyce używane są dwie wersje IP: IPv4 z adresami 32-bitowymi oraz IPv6
z adresami 128-bitowymi. IPv6 opracowano między innymi w odpowiedzi na
ograniczoną pulę adresów IPv4.

\subsection*{Adresowanie IPv4 i podział na podsieci}

Adres IPv4 składa się z 32 bitów zapisywanych zwykle jako cztery oktety
oddzielone kropkami. Oktet ma 8 bitów i może być przedstawiony w postaci
dziesiętnej z zakresu 0--255 albo jako zapis binarny złożony z ośmiu znaków.
Nie wszystkie adresy mogą być wykorzystywane jako zwykłe adresy hostów:
część zakresów ma specjalne zastosowania, a w typowej podsieci IPv4 osobne
adresy określają sieć i jej adres rozgłoszeniowy.

W pierwotnym modelu IPv4 adresy dzielono na klasy. W klasach A, B i C
domyślna część sieciowa miała odpowiednio 8, 16 i 24 bity. Klasa D służyła
do multicastu, czyli przesyłania pakietów do grupy odbiorców, a klasa E była
zarezerwowana. Podział klasowy jest obecnie historyczny -- współczesne sieci
korzystają z adresowania bezklasowego CIDR.

Jednym ze sposobów ograniczania zużycia publicznych adresów IPv4 jest
korzystanie z adresów prywatnych. Mogą one występować w wielu sieciach
lokalnych i w odróżnieniu od adresów publicznych nie są globalnie unikalne.
Zakresy prywatne opisane w RFC 1918 przedstawia tabela.

\begin{center}
\begin{tabularx}{\textwidth}{@{}lX@{}}
\toprule
Zakres adresów & Liczba adresów \\
\midrule
\texttt{10.0.0.0 -- 10.255.255.255} & 16\,777\,216 \\
\texttt{172.16.0.0 -- 172.31.255.255} & 1\,048\,576 \\
\texttt{192.168.0.0 -- 192.168.255.255} & 65\,536 \\
\bottomrule
\end{tabularx}
\end{center}

Maska podsieci pozwala wydzielić z adresu IP część określającą podsieć
i część określającą hosta. Ma taką samą długość jak adres: 32 bity dla IPv4
i 128 bitów dla IPv6. Zaczyna się ciągiem bitów o wartości 1, po którym
następuje ciąg bitów o wartości 0. Na podstawie adresu IP i maski można
wyznaczyć adres podsieci, adres rozgłoszeniowy oraz liczbę adresów hostów.
Maskę można zapisać binarnie, dziesiętnie z kropkami lub w notacji CIDR,
na przykład \texttt{255.255.240.0} odpowiada zapisowi \texttt{/20}.

VLSM (\textit{Variable-Length Subnet Mask}) umożliwia używanie masek
o różnych długościach, dzięki czemu rozmiar podsieci można dopasować do
jej potrzeb i efektywniej przydzielać adresy. CIDR
(\textit{Classless Inter-Domain Routing}) wykorzystuje zapis adresu wraz
z długością maski po ukośniku, na przykład \texttt{192.168.1.0/24}.
Ułatwia również agregowanie tras i ograniczanie rozmiaru tabel routingu.
VLSM jest wykorzystywany przy lokalnym planowaniu podsieci, natomiast CIDR
ma szersze zastosowanie w routingu.

\subsection*{Protokoły TCP i UDP}

TCP (\textit{Transmission Control Protocol}) jest połączeniowym protokołem
warstwy transportowej. Działa w modelu klient--serwer i zapewnia dostarczenie
pakietów we właściwej kolejności, bez duplikatów. Do nawiązania połączenia
wykorzystuje procedurę \textit{three-way handshake}:

\begin{enumerate}
    \item Klient wysyła segment SYN z początkowym numerem sekwencji.
    \item Serwer odpowiada segmentem SYN ze swoim numerem sekwencji oraz
    potwierdzeniem ACK otrzymania segmentu klienta.
    \item Klient wysyła ACK, potwierdzając otrzymanie segmentu SYN serwera.
\end{enumerate}

Połączeniem sterują flagi w nagłówku TCP. SYN inicjuje połączenie
i synchronizuje numery sekwencyjne, ACK wskazuje ważność numeru
potwierdzenia, FIN sygnalizuje zakończenie wysyłania danych, a RST resetuje
połączenie lub odrzuca próbę jego ustanowienia. Nagłówek TCP ma co najmniej
160 bitów i może zawierać opcje o zmiennej długości.

UDP (\textit{User Datagram Protocol}) jest bezpołączeniowym protokołem
warstwy transportowej. Nie gwarantuje dostarczenia wszystkich pakietów ani
zachowania ich kolejności, ale zapewnia szybszą transmisję i mniejszy
narzut danych niż TCP. Jego nagłówek składa się z czterech 16-bitowych pól:
portu źródłowego, portu docelowego, długości datagramu i sumy kontrolnej.

\subsection*{Porty}

Port jest mechanizmem programowym służącym do identyfikowania procesów
i usług sieciowych działających na urządzeniu. Razem z adresem IP tworzy
adres gniazda internetowego, wykorzystywany do przesyłania danych między
procesami na różnych urządzeniach. Port jest 16-bitową liczbą z zakresu
0--65535. IANA prowadzi rejestr portów przypisanych usługom, a zakres
podzielono na trzy grupy:

\begin{itemize}
    \item porty dobrze znane: 1--1023, używane przez wiele usług systemowych,
    między innymi FTP (20 i 21), DNS (53) oraz HTTP (80);
    \item porty zarejestrowane: 1024--49151, rezerwowane przez IANA na
    wniosek instytucji lub firmy;
    \item porty przydzielane dynamicznie: 49152--65535, przydzielane
    tymczasowo, między innymi na potrzeby usług i mechanizmów automatycznego
    przydzielania portów.
\end{itemize}

\subsection*{Skanowanie sieci za pomocą Nmap}

Nmap (\textit{Network Mapper}) jest skanerem portów TCP i UDP. Umożliwia
sprawdzanie dostępności hostów, rozpoznawanie usług oraz identyfikowanie
systemu operacyjnego skanowanego hosta. Rozpoznawanie systemu może opierać
się na charakterystycznych cechach pakietów, takich jak wartość TTL lub
numery sekwencyjne. Wśród rodzajów skanowania opisanych w instrukcji
znajdują się skanowanie TCP, SYN i UDP, wykrywanie aktywnych hostów,
rozpoznawanie wersji usług oraz systemu operacyjnego.

\end{document}